\section{The scale curve, and why its two metrics disagree}
\label{app:scale}

This appendix carries the characterization that \S\ref{sec:evaluation}
summarizes in a paragraph: what the three model sizes do to quality, whether
the closing of the gap to 4 bits itself slows, and why the answer depends on
which metric is asked. It is reported here rather than in the body because
none of it is a statement about the kernel or the format --- the transcoding
is a proved bijection of the decoded content, so the served path's quality
\emph{is} the artifact's --- but it is the context in which the 2-bit
operating point is or is not worth its kernel.

\textbf{Does the closing itself slow? On the MMLU gap, the bars decline to
say.} With the three gaps finally the same kind of number, the question the
word \emph{knee} asks can be put to a test, and on this metric the answer is
that one step is resolved and the other is not. Composing the two standard
errors in quadrature --- the campaigns are independent, and no pairing across
model sizes is available for a question set or would mean anything --- the gap
falls by 6.96 points from 4B to 8B ($\mathrm{SE}$ 1.82, $z = 3.82$,
$p = 0.0001$), by 1.40 from 8B to 14B ($\mathrm{SE}$ 1.68, $z = 0.83$,
$p = 0.40$), and by 8.36 end to end ($\mathrm{SE}$ 1.91,
$p \approx 10^{-5}$). The first closing is established and the total closing
is established; the second step is not. \emph{And $p = 0.40$ is not evidence
of constancy either.} On that step this metric is silent, and both readings
--- a curve that flattens, and a curve that keeps closing at a smaller pace
--- remain compatible with it. These are $z$-tests computed from the intervals
of Table~\ref{tab:scale}, not a further measurement. \emph{The same question
put to perplexity answers differently}, and it is answered below, once that
column has intervals of its own; a claim about the knee that does not name
its metric is half wrong whichever half it states.

\textbf{Perplexity gets error bars, and now at all three sizes.} Every
perplexity ratio in this paper was a point estimate while the MMLU column
carried intervals --- a comparison that tests its differences on one axis and
not the other invites exactly the question it should close. Nine paired
intervals now exist, three per model, formed window by
window over the 12 evaluation windows --- the same text in every arm, which is
what the matching fingerprint guarantees --- as a Student $t$ interval on the mean
per-window NLL difference, $n = 12$, 11 degrees of freedom, exponentiated to
a perplexity ratio. The exponentiation is exact rather than a delta-method
step, since every window scores the same token count. As excess over the
stated reference: at 4B, AWQ over FP16 $+10.49\%$ $[+8.55, +12.47]$, 2-bit
over FP16 $+38.45\%$ $[+33.62, +43.45]$, 2-bit over AWQ $+25.31\%$
$[+20.01, +30.84]$; at 8B, $+4.80\%$ $[+4.24, +5.35]$, $+22.01\%$
$[+19.37, +24.70]$ and $+16.42\%$
$[+14.17, +18.72]$; at 14B, $+3.81\%$ $[+3.27, +4.34]$, $+18.94\%$
$[+17.22, +20.68]$ and $+14.58\%$ $[+13.10, +16.08]$. None contains zero, and
all 108 window-level comparisons point the same way. The instrument returns a
null when it should: the same 8B file scored with an f16 and with an int8
embedding differs by $+0.0004\%$, interval $[-0.025, +0.026]$, which contains
zero.

\textbf{The 4B's three intervals rest on raw output, not on a summary.} A
campaign log that records aggregate perplexities and no per-window
log-likelihoods cannot be barred afterwards at any price short of rerunning
the arms, so the per-window lines are committed rather than cited by job
identifier (\texttt{docs/\allowbreak mesures/\allowbreak
a4-campagne-4b-ppl-BRUT-2026-08-06.txt}). Three controls ran before those
lines were used for anything: the three published perplexities replay from
the NLLs to the ten-thousandth (12.2369 / 13.5207 / 16.9422), the paired
differences are additive window by window --- all twelve residuals exactly
zero in double precision, which a one-window offset between two arms would
break and nothing else would show --- and the run fingerprint on all three
arms is the one the 8B and 14B campaigns carry.

One reservation bounds all nine intervals: they cover the sampling of the
evaluation corpus and nothing else, in particular not the calibration draw,
which is one seed per artifact and is unmeasured at these scales; they
therefore condition on the three files we produced. Both steps of the curve
are barred. The
8B-to-14B fall of the excess over FP16 is $-13.9\%$ with interval
$[-22.8, -4.9]$, and the 4B-to-8B fall is $-42.8\%$ with interval
$[-51.8, -33.5]$ --- the bar that the $-43\%$ quoted in this repository since
2026-08-10 had gone without. Against the AWQ reference --- the one the product
argument rests on --- the 8B-to-14B step reads $-1.58\%$ with interval
$[-3.14, -0.004]$,
$t = 2.2063$ against a threshold of 2.2010, and it clears zero by five
thousandths, so it must not be described as a significant closing. \emph{Those
two percentages are in different parameterizations and must not be read as a
ninefold difference}: $-13.9\%$ is the fall of the \emph{excess}, $-1.58\%$ the
fall of the \emph{ratio} of perplexities, which is the only form the journal
publishes for the AWQ reference. In the ratio parameterization the two
references read $-2.51\%$ $[-4.12, -0.88]$ and $-1.58\%$ $[-3.14, -0.004]$.
Data: \texttt{docs/\allowbreak mesures/\allowbreak
ppl-appariee-8b-14b-2026-08-17.txt} and
\texttt{docs/\allowbreak mesures/\allowbreak ppl-appariee-4b-2026-08-17.txt},
carried by \texttt{docs/data/ppl-appariee.csv}.

\textbf{On perplexity the closing does slow, and the knee resolves.} The
question put to the MMLU gap above can now be put to this
column, and here the three campaigns \emph{can} be paired across model sizes:
the 12 windows are the same text at every size, which the common fingerprint
guarantees, so the difference of differences is formed window by window rather
than by composing two standard errors. Against FP16 the perplexity
\emph{ratio} falls by a
factor $\times0.8812$ $[0.8561, 0.9071]$ from 4B to 8B, all 12 windows
agreeing in sign, and $\times0.9749$ $[0.9588, 0.9912]$ from 8B to 14B, 10 of
12 agreeing. \emph{Those are the same two steps as the $-42.8\%$ and
$-13.9\%$ of Table~\ref{tab:scale}, in the other parameterization}: 0.8812 is
the ratio of two perplexity ratios, $-42.8\%$ the fall of the excess, and the
warning two paragraphs above applies here too --- the test is run on the log
ratio because that is what the per-window differences are linear in, which is
what makes the pairing exact. The knee is the difference of those two log
ratios, and it is
$-0.1010$ $[-0.1377, -0.0643]$, $t = -6.06$ on 11 degrees of freedom, 11 of 12
windows agreeing: the first step lowers the excess strictly harder than the
second, with the near edge of the interval nearly four standard errors from
zero. Against the AWQ reference the same test still resolves but far less
comfortably --- $-0.0576$ $[-0.1011, -0.0140]$, $t = -2.91$ --- and it must not
be quoted with the confidence of the FP16 line. Data:
\texttt{docs/data/ppl-genou.csv}.

\textbf{Why the two metrics answer differently, and what neither buys.} They
are two quantities, not two measurements of one, and the most likely
explanation is power: perplexity averages over the 49\,140 scored tokens of
12 windows and pairs those windows across sizes, while MMLU is a rate over
2\,280 questions with no cross-size pairing available, so its test composes
two independent standard errors instead of cancelling a common nuisance. The
perplexity instrument is mechanically the sharper of the two. Two further
differences are real and neither is disposable. The columns do not measure the
same faculty --- 2 bits damages reasoning far more than recall and a
perplexity corpus mostly measures recall (\S\ref{sec:limitations}), so a curve
may bend on one and not the other. And they do not compare the same arms: the
perplexity knee is 2-bit against FP16, the MMLU knee is 2-bit against the
4-bit baseline. That last difference is checkable, and it does not carry the
divergence on its own --- against the AWQ reference the perplexity knee still
resolves. \emph{Hence every statement about the knee in this paper names its
metric}; an unqualified ``the knee holds'' is half wrong, and so is an
unqualified ``it does not''. What neither answer buys is the extrapolation.
The perplexity knee conditions on three artifacts with one calibration draw
each, which is the reservation above and bites hardest here since this test
conditions on all three at once; a resolved knee says the first step fell
harder than the second, not where the curve flattens. We still publish no
scaling law from three models, and a 32B point is what would settle the shape.
